\lysh{13318}{2}

\begin{alist}
\item $f(0) = -0 + \sqrt{2} = \sqrt{2}$.

$f(\sqrt{2}) = -\sqrt{2} + \sqrt{2} = 0$.

\[f(-\sqrt{2}) = -(-\sqrt{2}) + \sqrt{2} = \sqrt{2} + \sqrt{2} = 2\sqrt{2}.\]

\[[f(-\sqrt{2})]^2 = (2\sqrt{2})^2 = 2^2 \cdot (\sqrt{2})^2 = 4 \cdot 2 = 8.\]

\item Για $x = 0$, έχουμε βρει από το $\alpha)$ ερώτημα $y = f(0) = \sqrt{2}$, ενώ για $y = 0$ παίρνουμε
$0 = -x + \sqrt{2} \Leftrightarrow x = \sqrt{2}$, δηλαδή $f(\sqrt{2}) = 0$, τιμή που βρέθηκε στο $\alpha)$ ερώτημα.

Έτσι, βρήκαμε τα σημεία $B(0, \sqrt{2})$ και $A(\sqrt{2}, 0)$, τα οποία είναι αυτά στα οποία η γραφική παράσταση τέμνει τους άξονες $y'y$ και $x'x$ αντίστοιχα.
Παρατηρούμε ότι η συνάρτηση $f$ είναι της μορφής $y = f(x) = \alpha x + \beta$ με $x \in \mathbb{R}$, ($\alpha \cdot \beta \ne 0$)
οπότε η γραφική της παράσταση θα είναι μια ευθεία που δεν διέρχεται από την αρχή των
αξόνων και επομένως αρκεί να σχεδιάσουμε την ευθεία που διέρχεται από τα σημεία $A$ και $B$.
\end{alist}
